
\section*{Introduction}
Before delving into the theory of analytic rings and their associated notion of \emph{completeness} in chapter \ref{chapter: analytic rings and completeness}, we will in this chapter first recall the needed geometric foundations. In section \ref{section: schemes} we will recall the definition of spectra and schemes as ringed spaces that are locally such spectra. In section \ref{section: (quasi-)coherentness} we recall the notions of quasi-coherentness and coherentness of sheaves of modules. Section \ref{section: properness} recalls the notion of proper morphism between schemes as these are fundamental to duality. Then \ref{section: yoga} is concerned with introducing the basics of so-called \emph{$6$-functor formalisms} and their implications for duality in section \ref{section: coherent duality}.

\section{Schemes}
\label{section: schemes}

Recall the definition of affine spectra -- the basic building blocks of schemes.
\begin{definition}[Affine spectra, {
		\cite[\href{https://stacks.math.columbia.edu/tag/00DY}{Tag 00DY}]{stacks-project}
		\&
		\cite[\href{https://stacks.math.columbia.edu/tag/01HU}{Tag 01HU}]{stacks-project}}]
	\label{definition: affine spectra}
	
	Let $R$ be a commutative unitary ring. The \emph{(affine) spectrum of $R$} is the (locally) ringed space with underlying set $\Spec(R)$, the set of prime ideals of $R$. We endow this set with the Zariski topology, \ie closed sets are precisely those of the form
	$$V(I) \ldef \{P \in \Spec R \setseparator \forall f\in I\colon f\in P\}$$
	for some $I\subseteq R$. For each $f\in R$ call $D(f)\ldef \Spec R\setminus V(f)$ \emph{distinguished open}. These sets $D(f)$ form a basis of the Zariski topology, and the assignment $D(f)\mapsto R_f$ where $R_f$ is the localization of $R$ at $f$, forms a sheaf of rings on this basis. Hence, there is a sheaf of rings on $\Spec R$ denoted $\struct {\Spec R}$ extending this sheaf on the basis, \ie such that there is a natural isomorphism
	$$R_f \to \struct {\Spec R}(D(f))$$
	for any $f\in R$.
\end{definition}

%\begin{remark}[Interpretation of affine spectra]
%	Let $X=\Spec R$ be an affine spectrum. We think of $X$ as a geometric object which carries some additional structure, namely that we can specify on each open subset $U$ a collection $\struct X(U)$ of function-like objects on $U$. Since no prime ideal can contain a unit, we obtain $D(1)=X$ and hence $R\cong\struct X (D(1))=\struct X(X)$. So $R$ is the ring of \emph{functions} defined on all of $X$ and we call it the \emph{coordinate ring} of $X$. Compare this to the classical setting. Suppose we are given some affine complex variety $V$ with associated coordinate ring $\C[V] \ldef \{f\colon V\to \C\setseparator f\text{ is regular}\}$. As with affine schemes, the Zariski topology on $V$ admits a basis $D(f)=\{x\in V\setseparator f(x)\ne 0\}$ for the various $f\in \C[V]$ and the ring of regular functions on such a distinguished open is naturally isomorphic to the localization $\C[V]_f$ since expressions with denominator $f$ make sense on the non-vanishing set $D(f)$ of $f$. Hence the sheaf of regular functions $\struct V$ on $V$ will naturally identify with the structure sheaf $\struct {\Spec \C[V]}$ on $X$.
%	
%	Now since for a general affine scheme $X$ we should interpret elements $f\in \struct X(U)$ as regular functions on $U\subseteq X$, two questions arise: How do we evaluate such a function at a point $P\in U$ (which is by definition a prime ideal) and where should the values $f(P)$ live? To answer these we will restrict to the global case, which essentially corresponds to working with (globally defined) polynomials -- the general case follows formally similar but is a bit more contrived. Consider again an affine complex variety $V\subseteq \A^n$ and some polynomial function $f$ on $V$. Choose a point $p\in V$ with its corresponding maximal ideal $\mathfrak m < \C[V]$. There is a natural map of rings
%	$$\C[V] \to \C[V]/\mathfrak m \to \quot{\C[V]/\mathfrak m}\cong \C.$$
%	Under the isomorphism $\C[V] \cong \C[x_1,\dots, x_n]/I(V)$ the ideal $\mathfrak m$ corresponds to the ideal $(x_1-p_1,\dots, x_n - p_n)$ and the above map of rings corresponds to the evaluation $\C[x_1,\dots,x_n]/I(V)\to \C, f\mapsto f(p)$. Hence we can express the value $f(p)$ (which should correspond to $f(\mathfrak m)$ in the scheme theoretic case), by considering the image of $f$ along some suitable map of rings defined using the associated prime ideal!
%	
%	Indeed, coming back to the affine scheme $X=\Spec R$, with $f\in R$ and $P\in X$ a prime ideal, we define $f(P)$ to be the image of $f$ under the natural map
%	$$R\to R/P \to \quot{R/P}.$$
%	This answers both questions immediately, notice however that the ambient space of values obtained by evaluating at $P$, namely $\quot{R/P}$ varies with each point $P\in X$! This already hints at the fact these $f\in R$ should only be thought of a functions, but themselves are somehow something \emph{more}. For example such an $f$ might not even be determined by its values -- indeed for the affine spectrum $\Spec K[x]/(x^2)$ the global function $\overline x$ vanishes at every point, is however not the the same as the global function $\overline 0$ (simply by the fact that $\overline 0 \ne \overline x$ in the ring $K[x]/(x^2)$)! So some caution is needed. Nevertheless one can talk about the vanishing locus $\{X\in P\setseparator \forall f\in I\colon f(P)=0\}$ for any collection $I\subseteq R$ of functions. Since $f(P)=0$ if and only if $f\in P$ this recovers precisely the notion $V(I)$ from definition \ref{definition: affine spectra}.
%\end{remark}

For convenience, we define an affine scheme to be any locally ringed space that is isomorphic (as a locally ringed space) to an affine spectrum.
\begin{definition}[Affine schemes, {
		\cite[\href{https://stacks.math.columbia.edu/tag/01HW}{Tag 01HW}]{stacks-project}}]
	
	An \emph{affine scheme} is a locally ringed space $(X, \struct X)$ isomorphic to an affine spectrum as a locally ringed space.
\end{definition}

\begin{definition}[Schemes and their morphisms, {
		\cite[\href{https://stacks.math.columbia.edu/tag/01IJ}{Tag 01IJ}]{stacks-project}}]
	
	A \emph{scheme} is a locally ringed space $(X, \struct X)$ in which every point admits an open neighborhood $U$ such that $(U, \struct U)$ is an affine scheme. A morphism of schemes is a morphism of locally ringed spaces.
\end{definition}

One is usually only interested in morphisms of schemes whose fibers are not infinite-dimensional.
\begin{definition}[Being locally of finite type]
	\label{definition: being locally of finite type}
	
	A morphism $f\colon X\to Y$ of schemes is said to be \emph{locally of finite type} if for any point $x\in X$ there is an affine open neighborhood $\Spec A\cong U\subseteq Y$ of $x$ and an affine open neighborhood $\Spec R\cong V\subseteq Y$ of $f(x)$ such that $f(U)\subseteq V$ and such that the associated morphism of rings $R\to A$ induced by $f$ is \emph{of finite type}, \ie that $A$ is isomorphic as an $R$-algebra to a quotient of some polynomial algebra $R[X_1,\dots, X_n]$ over $R$.
\end{definition}

\begin{remark}[Interpreting being locally of finite type]
	In the setting of definition \ref{definition: being locally of finite type} if we write $A \cong R[X_1,\dots, X_n]/I$ for some ideal $I$ then we obtain that $U\cong \Spec A \cong \Spec R[X_1,\dots, X_n]/I \rdef V(I)$ is a closed subscheme of the affine space $\A_R^n$ over the space $\Spec R$. \Ie if $f\colon X\to Y$ is locally of finite type then it locally can be realized as a finite dimensional subscheme of affine space over (a part of) the base.
\end{remark}

There is an even stronger condition that essentially requires the fibers to be finite.
\begin{definition}[Being locally finite]
	\label{definition: being locally finite}
	
	A morphism $f\colon X\to Y$ of schemes is said to be \emph{locally finite} if for every point $x\in X$ there is an affine open neighborhood $\Spec A\cong U\subseteq Y$ of $x$ and an affine open neighborhood $\Spec R\cong V\subseteq Y$ of $f(x)$ such that $f(U)\subseteq V$ and such that the associated morphism of rings $R\to A$ induced by $f$ is \emph{finite}, \ie that $A$ is a finitely generated $R$-module.
\end{definition}

\begin{remark}[Interpreting finiteness]
	In the setting of definition \ref{definition: being locally finite} we find that since $R\to A$ is finite that by \cite[\href{https://stacks.math.columbia.edu/tag/05DR}{Tag 05DR}]{stacks-project} the fibers of $f\colon U\cong \Spec A\to \Spec R=V$ are finite, as desired. Thus, $f\colon X\to Y$ is best imagined as a branched cover of $Y$ (or rather of its image inside $Y$ -- the inclusion of a closed subset is for example finite as well, but many fibers are empty).
\end{remark}

A very well-behaved (for these we will obtain boundedness of the dualizing complex, see theorem \ref{theorem: exceptional direct and inverse image functors, the affine case}) class of (affine) schemes is given by \emph{complete intersections}.
\begin{definition}[Complete intersections]
	A morphism $f\colon \Spec A\to \Spec R$ of affine schemes is said the be a \emph{complete intersection} if there is a regular sequence $(f_1,\dots,f_r)$ in $R$ such that the natural morphism $R/\ideal{f_1,\dots, f_r}\to A$ of $R$-algebras induced by $R\to A$ is an isomorphism. Indeed, then $f$ is isomorphic (as schemes over $\Spec R$) to $V(f_1,\dots, f_r) = \Spec R/\ideal{f_1,\dots, f_r}$ into $\Spec R$.
\end{definition}


\section{(Quasi-)coherentness}
\label{section: (quasi-)coherentness}

Next we will introduce the notions of quasi-coherent and coherent sheaves on some scheme $X$. These essentially should correspond to sheaves of modules which locally look like modules respectively finitely generated modules. While this is not true in general (at least in the case of coherence) it is still a good way to think about them.

\begin{definition}[Quasi-coherent sheaves, {
		\cite[\href{https://stacks.math.columbia.edu/tag/01BE}{Tag 01BE}]{stacks-project}}]
	
	Let $(X, \struct X)$ be a ringed space. A sheaf $\sheaf F\in\structMod X$ is called \emph{quasi-coherent} if for every $x\in X$ there exists an open neighborhood $x\in U\subseteq X$ for which there is a presentation
	$$\struct U^{\oplus \vert J\vert} \to \struct U^{\oplus \vert I\vert} \to \sheaf F\restriction_U \to 0$$
	of $\sheaf F\restriction_U$, where $I$ and $J$ are possibly infinite sets, \ie if \emph{$\sheaf F$ is locally presentable}. We denote by $\qCoh X$ the full subcategory of quasi-coherent sheaves of $\structMod X$. If additionally for any $x$ the set $I$ and $J$ can be chosen finite, then $\sheaf F$ is called \emph{locally of finite presentation}.
\end{definition}

\begin{remark}[Terminology]
	Some authors (including the authors of The Stacks Project, \confer \cite[\href{https://stacks.math.columbia.edu/tag/01BM}{Tag 01BM}]{stacks-project}) call \emphquote{locally of finite presentation} just \emphquote{of finite presentation}. To the author this seems like a misnomer.
\end{remark}

\begin{definition}[Associated modules, {
		\cite[\href{https://stacks.math.columbia.edu/tag/01BH}{Tag 01BH}]{stacks-project}}]
	
	Let $(X, \struct X)$ be a ringed space and $\alpha\colon R\to \Gamma(X, \struct X)$ be any ring morphism. There is a functor $\associatedSheaf{(-)} \colon \mod R \to \structMod{X}$ left adjoint to $\Gamma(X, -)\colon \structMod{X} \to \mod R$ where for some sheaf $\sheaf F$ the $R$-module structure on $\Gamma(X,\sheaf F)$ is induced by $\alpha$. For an $R$-module $M$ the sheaf $\associatedSheaf{M}$ is called the \emph{sheaf associated to $M$ and $\alpha$}. If $R=\Gamma(X, \struct X)$ and $\alpha = \id_R$, then $\associatedSheaf{M}$ is simply called the \emph{sheaf associated to $M$}.
\end{definition}

\begin{definition}[Fundamental systems of neighborhoods]
	Let $X$ be a topological space and $x\in X$. A collection $\mathcal N\subseteq \mathcal N_x$ of neighborhoods of $x$ is called a \emph{fundamental system of neighborhoods of $x$} or \emph{neighborhood basis of $x$}, if for every neighborhood $N$ of $x$ there exists some neighborhood $N' \in \mathcal N$ with $N'\subseteq N$. Hence, if we define the partial order $\le$ by $N \le N' \Leftrightarrow N'\subseteq N$, then $\mathcal N$ is a fundamental system of neighborhoods for $x$ if and only if $\mathcal N$ is cofinal in $\mathcal N_x$.
\end{definition}

\begin{lemma}[Quasi-coherent sheaves are locally associated modules, {\cite[\href{https://stacks.math.columbia.edu/tag/01BK}{Tag 01BK}]{stacks-project}}]
	
	Let $(X,\struct X)$ be a ringed space and $x\in X$ a point with a fundamental system of quasi-compact neighborhoods. If $\sheaf F$ is a quasi-coherent $\struct X$-module, then there exists an open neighborhood $U$ of $x$ such that $\sheaf F\restriction_U$ is isomorphic to the sheaf $\associatedSheaf{M}$ associated to some $\Gamma(U, \struct X)$-module $M$.
\end{lemma}

For schemes this is always true.
\begin{corollary}[Quasi-coherent sheaves of schemes]
	Since in a scheme $X$ every point has a fundamental system of quasi-compact neighborhoods, namely the system of affine neighborhoods, we obtain that any quasi-coherent sheaf on $X$ is locally an associated module.
\end{corollary}

If the scheme in question is even affine, we obtain an even stronger statement. Quasi-coherent sheaves are nothing more than modules of the corresponding ring.
\begin{remark}[Quasi-coherent sheaves of affine schemes, {\cite[\href{https://stacks.math.columbia.edu/tag/01IB}{Tag 01IB}]{stacks-project}}]
	\label{remark: quasi-coherent sheaves of affine schemes}
	
	If $X=\Spec R$ is affine, then 
	$$\mod{R} \xleftrightarrows{\Gamma(X, -)}{\associatedSheaf{(-)}} \qCoh{\struct{X}}$$
	is an equivalence of categories.
\end{remark}

Thus, just like schemes are glued from affine schemes, quasi-coherent sheaves (on a scheme) are glued from modules associated to the local patches. A question arises: What about (quasi-coherent) sheaves that a locally \emph{finitely presented} (or equivalently \emph{finitely generated} if the scheme in question is locally Noetherian)? This will lead us to the notion of \emph{coherentness}. First we introduce the analogue of \emph{being finitely generated}.
\begin{definition}[Sheaves of finite type, {
		\cite[\href{https://stacks.math.columbia.edu/tag/01B5}{Tag 01B5}]{stacks-project}}]
	
	Let $(X, \struct X)$ be a ringed space. A sheaf $\sheaf F$ of $\struct X$-modules  is called \emph{of finite type} if for every $x\in X$ there is an open neighborhood $x\in U\subseteq X$ and a surjection
	$$\struct U^{\oplus n} \to \sheaf F\restriction_U$$
	for some $n\in \N$, \ie if $\sheaf F$ is locally generated by finitely many sections.
\end{definition}

Being \emph{coherent} is now given by requiring certain local maps to have finite type kernel. This is a priori stronger than requiring the sheaf to be locally of finite presentation as one needs many kernels to be of finite type. We will see in lemma \ref{lemma: coherentness vs being locally of finite presentation} that for a wide class of examples \emph{being coherent} and \emph{being locally of finite presentation} agree.
\begin{definition}[Coherent sheaves of modules, {
		\cite[\href{https://stacks.math.columbia.edu/tag/01BV}{Tag 01BV}]{stacks-project}}]
	
	Let $(X, \struct X)$ be a ringed space. A sheaf of $\struct X$-modules $\sheaf F$ is called \emph{coherent} if it is of finite type and for any open $U\subseteq X$ and any finite collection $s_1,\dots,s_n\in \sheaf F(U)$ the kernel of the corresponding morphism $\struct U^{\oplus n} \to \sheaf F\restriction_U$ is of finite type. The full subcategory of coherent sheaves of $\struct X$-modules is denoted $\Coh X$.
\end{definition}

A direct consequence is that any coherent sheaf is already quasi-coherent.
\begin{remark}[Coherent sheaves are quasi-coherent]
	Any coherent sheaf $\sheaf F$ on a locally ringed space $(X, \struct X)$ is locally of finite presentation and hence quasi-coherent. Indeed, since $\sheaf F$ is of finite type, choose for each $x\in X$ an open neighborhood $U$ and a surjection $\pi\colon\struct U^{\oplus n} \to \sheaf F\restriction_U$. By coherentness its kernel must be of finite type and thus itself admit a surjection $\struct V^{\oplus m} \to \kernel(\pi)\restriction_V$ for some $x\in V\subseteq U$. Hence, we obtain a local presentation
	$$\struct V^{\oplus m} \to \struct V^{\oplus n} \to \sheaf F\restriction_V\to 0$$
	showing that $\sheaf F$ is quasi-coherent.
\end{remark}

As mentioned before, if one imposes some regularity, then coherentness is nothing more than being locally of finite presentation.
\begin{lemma}[Coherentness vs being locally of finite presentation, {\cite[\href{https://stacks.math.columbia.edu/tag/01BZ}{Tag 01BZ}]{stacks-project}}]
	\label{lemma: coherentness vs being locally of finite presentation}
	
	Let $(X, \struct X)$ be a ringed space for which the structure sheaf $\struct X$ is coherent and let $\sheaf F$ be a sheaf of $\struct X$-modules. Then $\sheaf F$ is coherent if and only if it is locally of finite presentation.
\end{lemma}

\begin{remark}[For which schemes $X$ is $\struct X$ coherent?]
	The question now is: To which schemes $X$ does lemma \ref{lemma: coherentness vs being locally of finite presentation} apply? \Ie for which schemes $X$ is the structure sheaf $\struct{X}$ coherent. This is certainly true for any scheme that is locally Noetherian, \confer \cite[\href{https://stacks.math.columbia.edu/tag/01XZ}{Tag 01XZ}]{stacks-project}. This covers many examples of interest and certainly all varieties over fields. In general there is the notion of a coherent ring in the affine case (a ring $R$ whose associated module $\struct{\Spec R}$ is coherent, equivalently that all its finitely generated ideals are finitely presented).
\end{remark}

\section{Properness}
\label{section: properness}

Very important to the theory of duality is the notion of a proper morphism.

\begin{definition}[Closed immersions, {
		\cite[\href{https://stacks.math.columbia.edu/tag/01HK}{Tag 01HK}]{stacks-project}}]
	
	Let $(i, i^\sharp)\colon (Z,\struct Z)\to (X,\struct X)$ be a morphism of locally ringed spaces. We call $(i, i^\sharp)$ a \emph{closed immersion} if
	\begin{itemize}
		\item
		the map $i$ is a homeomorphism onto a closed subset of $X$,
		
		\item
		the morphism $i^\sharp\colon\struct X\to i_\ast \struct Z$ is an epimorphism, and
		
		\item 
		the $\struct X$-module $\sheaf I\ldef \ker i^\sharp$ is locally generated by sections.  
	\end{itemize}
\end{definition}

\begin{example}[Interpreting closed immersions]
	The most basic closed immersion is the case of the morphism of schemes $V(I)\ldef \Spec A/I \to \Spec A$ where $A$ is some ring and $I$ is some ideal of $A$. Every other closed immersion is built from these examples as by \cite[\href{https://stacks.math.columbia.edu/tag/01QO}{Tag 01QO}]{stacks-project} a morphism of schemes $(i, i^\sharp)\colon Z\to X$ is a closed immersion if and only if for each affine open $\Spec A\cong U \subseteq X$ the scheme $i^{-1}(U)$ over $\Spec A$ is isomorphic to $\Spec A/I \to \Spec A$ for some ideal $I$ of $A$, \ie $(i, i^\sharp)\colon Z\to X$ is a closed immersion if $Z$ inside $X$ is \emphquote{locally a vanishing set of some ideal}.
\end{example}

\begin{definition}[Quasi-compactness, {
		\cite[\href{https://stacks.math.columbia.edu/tag/01K3}{Tag 01K3}]{stacks-project} \& \cite[\href{https://stacks.math.columbia.edu/tag/005A}{Tag 005A}]{stacks-project}}]
	
	A morphism of schemes $f\colon X\to Y$ is said to be \emph{quasi-compact}, if the (set-theoretic) preimage $f^{-1}(V)$ of any compact open $V\subseteq Y$ is itself compact.
\end{definition}

\begin{definition}[(Quasi-)separatedness, {
		\cite[\href{https://stacks.math.columbia.edu/tag/01KK}{Tag 01KK}]{stacks-project}}]
	Let $f\colon X\to Y$ be a morphism of schemes and consider the diagonal morphism $\Delta\colon X\to X\times_Y X$. If $\Delta$ is a closed immersion, then $f$ is called \emph{separated}. If $\Delta$ is quasi-compact then $f$ is called \emph{quasi-separated}. By \cite[\href{https://stacks.math.columbia.edu/tag/01K7}{Tag 01K7}]{stacks-project} any separated morphism is also quasi-separated.
\end{definition}

\begin{remark}
	Imposing quasi-separatedness ensures that geometrically pathological cases like the affine line with doubled origin $(\A_K^1\sqcup \A_K^1)/(\A_K^1\setminus\{0\}) \to \Spec K$ is excluded from consideration. Including such examples would be an obstruction to developing a good local theory as in non-(quasi-)separated spaces weird local artifacts occur: The two origins are distinct, yet there are no two \quote{essentially} disjoint closed subsets separating these two points. In an ordinary variety (over an algebraically closed field) one would expect for any two given distinct points to find two (up to a codimension $1$ subvariety) disjoint curves, one passing through the first point one passing through the second one. This is not the case in the affine line with doubled origin. Even worse, there are two ways to complete the inclusion of the smooth locally closed subvariety $\A_K^1\setminus\{0\}$ to a closed smooth subvariety. We essentially loose \quote{uniqueness of limits}. 
\end{remark}

%\begin{example}
%	Let $K$ be a field, $n\ge 2$ and $X$ be the affine $n$-space with two origins, \ie the scheme obtained from gluing two copies $X_1$ and $X_2$ of $\A_K^n\ldef \Spec K[x_1,\dots, x_n]$ along $\A_K^n\setminus \{0\}$ with $0\ldef\langle x_1,\dots,x_n\rangle$. Then $X\to \Spec K$ is not quasi-separated since the diagonal $\Delta\colon X\to X\times_k X$ is not quasi-compact. Indeed the inverse image of the affine $X_1\times_k X_2 \subseteq X\times_k X$ is $X_1\cap X_2\cong \A_K^n\setminus\{0\}$ which is neither affine nor compact by \cite[\href{https://stacks.math.columbia.edu/tag/01IL}{Tag 01IL}]{stacks-project} since $n\ge 2$.
%\end{example}

\begin{definition}[Universally closed morphisms]
	Let $f\colon X\to Y$ be a morphism of schemes. Then $f$ is called \emph{universally closed} if the map of underlying topological spaces is closed and any base change $X\times_Y Z\to Z$ of $f$ along a morphism $h\colon Z\to Y$ of schemes is closed as well.
\end{definition}

\begin{definition}[Properness, {
		\cite[\href{https://stacks.math.columbia.edu/tag/01W1}{Tag 01W1}]{stacks-project}}]
	
	A morphism $f\colon X\to Y$ of schemes is called \emph{proper} if it is separated, universally closed and of finite type.
\end{definition}


\begin{example}[Interpretation of properness]
	There is an alternative characterization of properness that is quite geometrically pleasing. Suppose that $f\colon X\to Y$ is a quasi-separated morphism of finite type. Then by \cite[\href{https://stacks.math.columbia.edu/tag/0BX4}{Tag 0BX4}]{stacks-project} the morphism $f$ is proper if and only if for every valuation ring $A$ with field of fractions $K$ fitting into a commutative solid diagram
	$$\begin{tikzcd}
		\Spec K \ar[r]\ar[d] &X\ar[d]\\
		\Spec A\ar[r]\ar[ur, dotted] &Y
	\end{tikzcd}$$
	there is a unique dotted arrow making the diagram commute. Observe that $\Spec K$ is the generic point of $\Spec A$ and that $\Spec A\to Y$ can be imagined as a parameterized (by the one-dimensional space $\Spec A$) smooth curve in $Y$. Each such solid diagram then essentially provides a lift of the generic point (so of \quote{almost all} of the curve) to $X$. The existence and uniqueness of the dotted arrow then assert that the missing points in the lifted curve can be uniquely imputed to obtain a lift of the curve in $Y$ to a smooth curve in $X$. So properness is a kind of \emph{(relative) completeness} for schemes.
	
	This for example is clear in case of a closed immersion $X\hookrightarrow Y$: If \quote{almost all} of a curve $S\subseteq Y$ already lies in $X$ then so does the whole curve: $X$ is closed in $Y$ so already contains all its limiting points. Properness only gets properly interesting in cases where the underlying map is surjective. \Eg the projection $\A_\C^2\to \A_\C^1$ to the first component fails to be proper: Denote by $x$ and $y$ the coordinate functions of the coordinate ring of $\A^2_\C$. The \quote{almost} lift of the curve $\A^1_\C\to \A^1_\C, t\mapsto t$ to the curve $\A^1_\C \to \A^2_\C, t\mapsto (t, 1/t)$ with image in $V(xy - 1)$ already fails to lift. This of course is the starting point for the cherished concept of a projective space.
\end{example}


Proper morphisms of schemes are both quasi-separated (by definition) and quasi-compact. Indeed, a proper morphism is universally closed by definition and hence quasi-compact by \cite[\href{https://stacks.math.columbia.edu/tag/04XU}{Tag 04XU}]{stacks-project}. Along qcqs morphisms it is possible to transport quasi-coherent sheaves.

\begin{lemma}[Direct and inverse images of quasi-coherent sheaves]
	\label{lemma: direct and inverse images of quasi-coherent sheaves}
	
	Let $f\colon X\to Y$ be a morphism of ringed spaces. If $\sheaf G \in \qCoh{Y}$, then also $f^\ast \sheaf G \in \qCoh X$. If $f$ is a qcqs morphism of schemes and $\sheaf F\in\qCoh X$, then also $f_\ast\sheaf F \in \qCoh Y$. Hence, if $f$ is qcqs morphism of schemes, direct and inverse image restrict to an adjoint pair
	$$\qCoh X \xleftrightarrows{f^\ast}{f_\ast} \qCoh Y.$$
\end{lemma}
\begin{proof}
	For inverse images this is \cite[\href{https://stacks.math.columbia.edu/tag/01BG}{Tag 01BG}]{stacks-project}, for direct images this is \cite[\href{https://stacks.math.columbia.edu/tag/01LC}{Tag 01LC}]{stacks-project}.
\end{proof}

\begin{remark}
	As any morphism$ f\colon \Spec B \to \Spec A$ of affine schemes is qcqs by \cite[\href{https://stacks.math.columbia.edu/tag/01S7}{Tag 01S7}]{stacks-project}, the previous lemma \ref{lemma: direct and inverse images of quasi-coherent sheaves} directly applies. This is not very surprising 
	because if $f$ corresponds to $\phi \colon A\to B$ then restricted to quasi-coherent modules, $f_\ast$ and $f^\ast$ correspond to scalar restriction and scalar extension of modules along $\phi$.
\end{remark}

\section{Yoga}
\label{section: yoga}

The natural starting point for coherent duality (and many related duality theories) is the idea of a \emph{$6$-functor formalism}. In general, a $6$-functor formalism ascribes to each object of a category of geometric objects $C$ (be it schemes or some category of locally compact topological spaces) a category (best thought of as a derived category of sheaves) and then relates them by $6$ classes of functors -- two of which (the \emphquote{exceptional} ones) are only defined on a pre-specified subclass $E$ of morphisms of $C$. Manipulating these $6$-functors is also called \emph{the yoga of $6$-functors}.

There is a rich theory surrounding this very general formalism -- this is for example covered in Scholze's notes \cite{sixfunctors} -- but we are only interested in one specific application (namely coherent duality). Thus, let us not introduce the general theory and instead consider by hand what this formalism should be in our case.

%Indeed tensor products and sheaf-$\Hom$s of coherent sheaves are again coherent by \cite[\href{https://stacks.math.columbia.edu/tag/01CE}{Tag 01CE}]{stacks-project} and \cite[\href{https://stacks.math.columbia.edu/tag/01CQ}{Tag 01CQ}]{stacks-project} respectively, so the two functors restrict to coherent sheaves.

For us, the category of geometric objects is of course a subcategory $X$ of $\Schemes$, the category of schemes. Classically one attaches to each scheme $X$ a derived category $\derived{X}$ of sheaves of $\struct X$-modules, \eg coherent or quasi-coherent sheaves. Let us postpone the discussion of the appropriate $E$ for now and let us specify instead the $6$ classes of functors. The first $2$ classes are given by
\begin{enumerate}
	\item
	the (derived) tensor products $-\otimes^\LeftD_{\struct X} -$
	
	\item
	and their (partial) right adjoints $\intRHom_{\struct X}$ the (derived) sheaf of homomorphisms,
\end{enumerate}
making each $\derived{X}$ closed symmetric monoidal. The next $2$ classes of functors associate to any morphism $f\colon X\to Y$ of schemes the
\begin{enumerate}
	\item[3.]
	(derived) direct image $f_\ast\colon \derived{X}\to \derived{Y}$ of $\struct X$-modules and
	
	\item[4.]
	its left adjoint $f^\ast\colon \derived{Y}\to \derived{X}$, the (derived) inverse image of $\struct Y$-modules.
\end{enumerate}
These are the derived functors of the usual direct and inverse image $f_\ast$ and $f^\ast$ and hence should maybe be more appropriately notated by $\RightD f_\ast$ and $\LeftD f^\ast$. It is often however more convenient to conflate the names and the symbols used to describe these functors. In particular, for a general $6$-functor formalism these do not even have to be derived functors nor do the categories $\derived{X}$ need to be derived (in Scholze's \cite{sixfunctors} they are even \quote{generic} $\infty$-categories). The final two classes associate to any morphism $f\colon X\to Y$ of $E$ the
\begin{enumerate}
	\item[5.]
	\emph{exceptional direct image} $f_!\colon \derived{X} \to \derived{Y}$ and
	
	\item[6.]
	its right-adjoint $f^!\colon \derived{Y} \to \derived{X}$, the \emph{exceptional inverse image}.
\end{enumerate}
These $6$ classes of functors now are subjected to several constraints and compatibilities they ought to satisfy. One such constraint is that one requires that $f_!$ agrees with $f_\ast$ in case $f$ is proper. Another example is the \emph{projection formula}: If $f\colon X\to Y$ is proper then for all $\sheaf F\in\derived{X}$ and $\sheaf G\in\derived{Y}$ there is a natural isomorphism
$$f_\ast\sheaf F\otimes_\struct{Y} \sheaf G \to f_\ast(A\otimes_{\struct{X}} f^\ast \sheaf G)$$
in $\derived{Y}$.

Classically there area few choices. If one wants to restrict to (quasi-)coherent sheaves then one needs to be restrictive in the definition of $C$ to ensure that for any morphism $f$ in $C$ the two functors $f_\ast$ and $f^\ast$ are well-defined, see lemma \ref{lemma: direct and inverse images of quasi-coherent sheaves}. Further, to find a large class $E$ usually requires some thought, as it is hard to construct an appropriate functor $f_!$. In case of sheaves on the étale site of a scheme, one obtains such a functor $f_!$ as the \emph{direct image with compact support}, see \cite[\href{https://stacks.math.columbia.edu/tag/0F4W}{Tag 0F4W}]{stacks-project}. In general, it is not so clear how to make a good choice of $f_!$ (and hence of $f^!$).

For us this issue will be resolved using condensed mathematics. Clausen and Scholze introduce in \cite{condenseddotpdf} for every schemes $X$ a category $\derived{\struct{X,\blacksquare}}$ and for a large class of morphisms an exceptional direct image $f_!$. In the affine case of $X=\Spec A$ this category $\derived{\struct{X, \blacksquare}}$ will be the derived category $\derived{A_\blacksquare}$ of so-called \emph{$A_\blacksquare$-complete} condensed $A$-modules (see chapter \ref{chapter: analytic rings and completeness}) while the $f_!$ will be defined for any $f\colon \Spec A\to \Spec R$ where $R$ is Noetherian and $f$ is of finite Tor-dimension (see chapter \ref{chapter: affine duality}). These definitions make critical use of condensed mathematics and provide a new approach to constructing a good derived theory of schemes.

%\begin{todo}
%	More on $\compactlySupportedCohomology$
%	
%	Update the conditions for the existence of $f_!$.
%\end{todo}
%
%Inspired by étale-cohomology one gives a special name to the cohomology of $f_!$. If $f\colon X\to \Spec K$ is a scheme over the field $K$ such that $f_!$ is defined then one also writes $\RightD\compactlySupportedGlobalSections(-, X) \ldef f_!$ and calls (either the total functor or its cohomology groups $\compactlySupportedCohomology^n(X, -)\ldef \RightD^n\compactlySupportedGlobalSections(-, X)=H^n\circ f_!$) the \emph{compactly supported cohomology of $X$}. If $f$ is proper then $f_!=\RightD f_\ast$ and compactly supported cohomology agrees with ordinary cohomology.


\section{Coherent Duality}
\label{section: coherent duality}

One of the ways to formulate coherent duality as follows.
\begin{theorem}[Coherent Duality, {\cite[Chapter VII, Corollary 3.4 (a) \& (c)]{residues}}]
	\label{theorem: coherent duality}
	
	Let $f\colon X\to Y$ be a proper morphism of finite type between \emph{Noetherian} (locally Noetherian and quasi-compact) schemes of finite Krull dimension admitting a dualizing complex. Then the right adjoint $f^!\colon \derivedPlus{\Coh{Y}} \to \derivedPlus{\Coh{X}}$ of $\RightD f_\ast\colon  \derivedPlus{\Coh{X}} \to \derivedPlus{\Coh{Y}}$ exists and the natural morphism
	$$\RightD f_\ast \RightD\intHom_{\struct X}^\bullet(F^\bullet, f^! G^\bullet) \to \RightD\intHom_{\struct Y}^\bullet(\RightD f_\ast F^\bullet, G^\bullet)$$
	is a (quasi)isomorphism for all complexes $F^\bullet \in \derivedMinus{\qCoh{X}}$ and $G^\bullet \in \derivedPlus{\Coh{Y}}$.
\end{theorem}

\begin{remark}[Existence of a dualizing complexes, {\cite[V \S10, sufficient conditions]{residues}}]
	When does a dualizing complex exist?
	\begin{itemize}
		\item
		If the scheme $X$ is Gorenstein (see \cite[\href{https://stacks.math.columbia.edu/tag/0AWV}{Tag 0AWV}]{stacks-project}) and of finite Krull dimension then $\struct{X}[0]$ is a dualizing complex for $X$. This covers the case of $X$ being Cohen--Macaulay (see \cite[\href{https://stacks.math.columbia.edu/tag/02IN}{Tag 02IN}]{stacks-project}) and of course the case where $X$ is regular and locally Noetherian (the definition of a Gorenstein scheme and of a Cohen--Macaulay scheme require local Noetherianness).
		
		\item
		If $f\colon X\to Y$ is a morphism of finite type with $Y$ Noetherian then if $Y$ admits a dualizing complex $\canonical_Y^\bullet$ so does $X$. If $f$ is proper with right adjoint $f^!$ then this is $\canonical_X^\bullet \ldef f^!\canonical_Y^\bullet$. This of course covers the case where $Y=\Spec K$ for some field $K$ as $Y$ is regular and (locally) Noetherian.
	\end{itemize}
\end{remark}

In the case where $f\colon X\to \Spec K$ is proper over the field $K$ this can be simplified a bit. Notice that if $f$ is of finite type with $X$ quasi-compact the Noetherianness is automatic.
\begin{corollary}[Coherent duality, a special case]
	\label{corollary: coherent duality, a special case}
	
	Let $f\colon X\to \Spec K$ be a morphism of finite type such that $X$ is quasi-compact and of finite Krull dimension. Denote by $(-)^\vee$ the $K$-dual. Then in $\derived{K}$ there is an isomorphism
	$$\eRHom_\struct{X}(F^\bullet, f^!K)\cong \RightD\Gamma(F^\bullet, X)^\vee$$
	%		$$\eHom_\derived{\struct{X}}(F^\bullet, f^!K) \cong \eHom_\derived{\struct{X}}(\struct{X}, F^\bullet)^\vee$$
	natural in the complexes $F^\bullet \in \derivedMinus{\qCoh{X}}$.
\end{corollary}
\begin{proof}
	First, observe that $f_\ast$ is nothing more than the global sections functor $\Gamma(X, -)\cong \eHom_\struct{X}(\struct{X}, -)$. Furthermore, $\Gamma(X, -) \circ \intHom_\struct{X} = \eHom_\struct{X}$ gives that $\RightD f_\ast\intRHom_\struct{X} \cong \eRHom_\struct{X}$. As $(-)^\vee$ is short for $\eRHom_K(-, K)$ we find that theorem \ref{theorem: coherent duality} implies
	$$\eRHom_\struct{X}(F^\bullet, f^!K)\cong \RightD\Gamma(F^\bullet, X)^\vee$$
	for $G^\bullet = K$.
\end{proof}
This of course includes all quasi-compact \emph{varieties} (reduced, separated schemes of finite type over a(n algebraically closed) field).


\begin{remark}
	\label{remark: upper shriek of structure sheaf is concentrated}
	
	Suppose we are again in the setting of theorem \ref{theorem: coherent duality} and recall the definition of a Cohen--Macaulay morphism of schemes from \cite[\href{https://stacks.math.columbia.edu/tag/045Q}{Tag 045Q}]{stacks-project}.	In case that $f\colon X\to \Spec Y$ is Cohen--Macaulay (this includes the case where $f$ is smooth) such that all fibers have the same dimension $d\in\N$ we find that $f^!\struct{Y}$ is a coherent sheaf concentrated in degree $d$ by \cite[Chapter VII, \S4]{residues}.
\end{remark}

One immediately obtains the following corollary of coherent duality. This is the version of Serre duality presented in \cite[Chap. III, Thm. 7.6]{algebraic-geometry} for projective schemes.
\begin{corollary}[Serre duality for Cohen--Macaulay schemes]
	\label{corollary: serre duality for cohen--macaulay schemes}
	
	If $X\to \Spec K$ is a finite type proper Cohen--Macaulay scheme over some field $K$ of pure dimension $n\in \N$ then there is a coherent sheaf $\canonical_{X/K}$ such that for all $k\in \Z$ there is an isomorphism
		$$\Ext_{\struct X}^k(\sheaf F, \canonical_{X/K}) \iso H^{n-k}(X, \sheaf F)^\dual$$
natural in the coherent sheaf $\sheaf F$.
\end{corollary}
\begin{proof}
Observe that by remark \ref{remark: upper shriek of structure sheaf is concentrated} the complex $f^!K$ is a coherent sheaf concentrated in degree $n$. Denote by $\canonical_{X/K}$ this sheaf such that $f^!K = \canonical_{X/K}[n]$. By corollary \ref{corollary: coherent duality, a special case} we find that
$$\eRHom_\struct{X}(\sheaf F, \canonical_{X/K}[n]) \cong \RightD\Gamma(\sheaf F, X)^\vee$$
and application of $H^{k-n}=H^k\circ (-)[-n]$ yields that
$$H^k\eRHom_\struct{X}(\sheaf{F}, \canonical_{X/K}) \cong (H^{n-k}\RightD\Gamma(\sheaf F, X))^\vee$$
since $\eRHom_\struct{X}$ is compatible with shifts and since $(-)^\vee$ is exact, however reverses the complexes. Thus,
$$\Ext^k_\struct{X}(\sheaf F, \canonical_{X/K})\cong H^{n-k}(X, \sheaf{F})^\vee.$$
\end{proof}

%\begin{theorem}[Serre Duality for Cohen-Macaulay schemes]
%	\label{theorem: Serre duality for Cohen-Macaualay schemes}
%	If $X$ is a proper Cohen-Macaulay scheme over some field $K$ of relative dimension $n$, then there exists a coherent sheaf $\canonical_X$ such that there are natural isomorphisms
%	$$\Ext_{\struct X}^k(\sheaf F, \canonical_X) \iso H^{n-k}(X, \sheaf F)^\dual$$
%	for all $k\in \Z$, natural in the coherent sheaf $\sheaf F$. If $i\colon X\to Y$ is a closed immersion into some smooth scheme $Y$ proper over $K$ such that $X$ is of codimension $r$ in $Y$, then $\canonical_X = i^\ast\intExt_{\struct Y}^r(i_\ast\struct X, \canonical_Y)$. If additionally $X$ is a complete intersection then $\canonical_X = \canonical_Y\restriction_X \structTensor{X} \exterior^r N_{X/Y}$ where $N_{X/Y}$ is the normal bundle of $X$ in $Y$, a locally free $\struct X$-module of rank $r$.
%\end{theorem}
In case that the coherent sheaf $\sheaf F$ in corollary \ref{corollary: serre duality for cohen--macaulay schemes} is locally free, one rephrase the obtained isomorphism in more familiar terms. For this we need the following characterization of internal $\Hom$s out of a locally free sheaf.
\begin{lemma}
\label{lemma: sheaf hom for vector bundles}

Let $X$ be a scheme. If $\sheaf L$ is locally free of finite rank and $\sheaf G$ is any sheaf of $\struct X$-modules, then
$$\intHom_{\struct X}(\sheaf L, \sheaf G) \iso \sheaf L^\dual \structTensor{X} \sheaf G.$$
If $X$ is a scheme over $\Spec K$ for some field $K$, then this isomorphism is $K$-linear.
\end{lemma}
\begin{shortproof}
This can be calculated locally where it is clear due to freeness.
\end{shortproof}

\begin{corollary}
\label{corollary: ext of vector bundle}

In the setting of the previous lemma \ref{lemma: sheaf hom for vector bundles} we obtain that
$$\Ext^\bullet_{\struct X} (\sheaf L, -) \ldef \eStructRHom{X}(\sheaf L, -) \iso \RightD \Gamma(X, \sheaf L^\dual \structTensor{X} -) \rdef H^\bullet(X, \sheaf L^\dual \structTensor{X} -)$$
as functors $\derived{\struct{X}}\to\derived{\Ab}$. If $X$ is a scheme over $\Spec K$ for some field $K$ then this is an isomorphism in $\derived{K}$ as well.
\end{corollary}
\begin{proof}
We have
\begin{align*}
	\eStructHom{X}(\sheaf L, -)&\iso \eStructHom{X}(\sheaf L\structTensor{X} \struct X, -)\\
	&\iso \eStructHom{X}(\struct X, \intHom_{\struct X}(\sheaf L, -))\\
	& \iso \eStructHom{X}(\struct X, \sheaf L^\dual\structTensor{X} -)\\
	&\iso \Gamma(X, \sheaf L^\dual \structTensor{X} -).
\end{align*}
If $X$ is a scheme over $\Spec K$ then $\mod{\struct{X}}$ is $K$-enriched and these isomorphisms are all of $K$-modules. It follows that the derived functors
$$\eStructRHom{X}(\sheaf L, -) \iso \RightD\Gamma(X, \sheaf L^\dual\structTensor{X} -)$$
agree as well as $\sheaf{L}^\vee\otimes_\struct{X}-$ is exact.
\end{proof}

We now can reformulate Serre duality in its classical form.
\begin{corollary}[Serre Duality]
\label{theorem: serre duality}

	If $X\to \Spec K$ is a finite type proper Cohen--Macaulay scheme over some field $K$ of pure dimension $n\in \N$ then there is a coherent sheaf $\canonical_{X/K}$ such that for all $k\in \Z$ there is an isomorphism
	$$H^k(X, \sheaf F) \cong H^{n-k}(X, \sheaf F^\vee \otimes_{\struct{X}} \canonical_{X/K})^\dual$$
	natural in the locally free sheaf $\sheaf F$ of finite rank.
	%	Let $X\to \Spec K$ be a proper, smooth variety over a field $K$ of (relative) dimension $n$. Denote by $\canonical_X\ldef \exterior^n T^\ast X$ the sheaf of $n$-forms on $X$. Then for any locally free sheaf $\sheaf F$ of finite rank we obtain
	%	$$H^k(X, \sheaf F) \iso H^{n-k}(X, \sheaf F^\dual \structTensor{X} \canonical_X)^\dual$$
	%	for all $k\in \Z$.
\end{corollary}
\begin{proof}
	This follows by noticing that $\Ext^k_{\struct{X}}(\sheaf F, \canonical_{X/K})\cong H^k(X, \sheaf F^\vee \otimes_{\struct{X}} \canonical_{X/K})$, by applying the (self-inverse) $(-)^\vee$ and replacing $k$ by $n-k$.
\end{proof}


%\section{Basic definitions and Coherent Duality}
%
%\begin{definition}[Support of a section, {
%		\cite[\href{https://stacks.math.columbia.edu/tag/01AT}{Tag 01AT}]{stacks-project}}]
%define this for any ringed space?
%	Let $X$ be a scheme and $\sheaf F$ a sheaf of $\struct X$-modules. If $U$ is an open set of $X$ and $s\in\sheaf F(U)$ is a section, then the \emph{support of $s$} is the set $\support{s}\ldef \{x\in U \setseparator 0\ne s_x\in\sheaf F_x\}$ of points of $U$ where the stalk of $s$ does not vanish. By \cite[\href{https://stacks.math.columbia.edu/tag/01AU}{Tag 01AU}]{stacks-project}, the support $\support s$ is closed in $U$.
%\end{definition}
%
%\begin{definition}[Exceptional/proper direct image]
%	Let $f\colon X\to Y$ be a (separated, finite type?) morphism of schemes. If $\sheaf F$ is a sheaf on $X$, $V$ is an open set of $Y$ and $s\in (f_\ast \sheaf F)(V)=\sheaf F(f^{-1}(V))$ is a section, then $\support s$ is a closed subset of the open $f^{-1}(V)$. Thus the support naturally carries the structure of a scheme and $f$ restricts to a morphism $f\restriction_{\support s}\colon \support{s}\to V$. Call the section $s$ \emph{of proper support}, if $f\restriction_{\support s}$ is a proper morphism of schemes. With this in mind, define a functor
%	$$f_!\colon \mod{\struct X}\to \mod{\struct Y}$$
%	such that for any sheaf $\sheaf F\in\mod{\struct X}$ and any open set $V\subseteq Y$
%	$$(f_! \sheaf F)(V) \ldef \{s\in\sheaf (f_\ast\sheaf F)(V)\setseparator\text{$s$ is of is proper support}\}.$$
%	Then $f_!$ is called the \emph{exceptional or proper direct image}. The natural inclusion $f_!\sheaf F\hookrightarrow f_\ast \sheaf F$, turns $f_!$ into a subfunctor of $f_\ast$.
%\end{definition}
%
%If $X$ is an $k$-scheme for some field $k$ and $f\colon X\to \Spec k$ is the projection, then the direct image $f_\ast$ naturally identifies with the global sections functor $\Gamma(X, -)$. Under this correspondence, $f_!$ identifies with the functor $\compactlySupportedGlobalSections(X, -)$, assigning to each sheaf of $\struct X$-modules its global sections of proper support. This leads to the following
%\begin{definition}[Compactly supported cohomology]
% Is it really left-exact? Maybe we need separated and finite type for that
%	Let $X$ be a $k$-scheme and $f\colon X\to\Spec k$ the projection. Then the functor $\compactlySupportedGlobalSections(X, -)\colon \mod{\struct X}\to \mod{k}$ is called the \emph{global sections with compact support}. It is left exact and has a right derived functor $\RightD\compactlySupportedGlobalSections(X,-)$. The cohomology groups $\compactlySupportedCohomology^\bullet(X, \sheaf F)\ldef \RightD^\bullet\Gamma_c(X, \sheaf F)$ are called the \emph{cohomology with compact support of $X$ with values in the sheaf $\sheaf F$}. The natural map $\compactlySupportedGlobalSections(X, -)\to \Gamma(X, -)$ induces is a comparison map $\compactlySupportedCohomology^\bullet(X, -)\to H^\bullet(X, -)$.
%\end{definition}
%
%\begin{remark}
%	If $f\colon X\to Y$ is (separated, of finite type and) proper, then every section is of proper support and thus $f_\ast = f_!$. In particular $\compactlySupportedGlobalSections(X, -)=\Gamma(X, -)$ and the comparison map $\compactlySupportedCohomology^\bullet(x, -) \to H^\bullet(X, -)$ is an isomorphism.
%\end{remark}

%\begin{remark}
%	Let $f\colon X=\Spec A\to S=\Spec R$ be a morphism of affine, finite type $\Z$-schemes. Define $H_!^\bullet(X, -) \ldef H^\bullet f_!$. Theorem \ref{theorem: existience of upper shriek for affines} asserts, that $f_!$ preserves compact objects. A discrete compact object of $D(A_\blacksquare)$
%\end{remark}
